\chapter{Purpose and scope}
\label{ch:purpose}

\section{What Blendity is}
Blendity is a 3D editor written in C++17. Its windows, navigation, shortcuts, gizmos, Inspector and Play mode follow
the Unity editor \cite{UnityManual2026}; underneath, its meshes, Edit Mode tools, UV unwrapping, Principled materials,
render engines (an EEVEE-like rasterizer and a Cycles-like path tracer on the CPU and any Vulkan GPU), skins and IK
follow Blender's design and source code \cite{Blender2026}. It is not a fork: each subsystem is re-implemented, following
Blender's design, so the editor builds and runs with nothing but a compiler, and the matching Blender source file is
listed for every module. When Blender's prebuilt libraries are present it links the same ones Blender uses, each behind
a switch with Blendity's own code as the fallback (\adr{0001}).

It is one process with no plugins and no scripting runtime. Three executables share one static library built from
\code{src/}: \code{Blendity} (the editor), \code{blendity\_tests} and \code{blendity\_stress}. Everything a tool does is
reachable from the Console's command line, which is how the tests, the headless mode and the HCI harnesses drive it.

\section{Who it is for}
The requirements (\code{docs/requirements.md}) name the audience by what they must be able to do: a Unity-style editor
(R-01) with Blender-equivalent modeling (R-02) and SketchUp / UModeler / Plasticity-style direct modeling (R-03), UVs,
materials and textures (R-04), rendering with both engines (R-05), file formats matching Blender 5.2 (R-06), every tool
reachable from menus, panels, rebindable shortcuts and the console (R-07), adjustable operations (R-08), a Learn tab
(R-09), camera filters (R-10), Unity-like global illumination (R-11, R-12), physics for animation as in Cascadeur
(R-13), fingers and a pen on a 2-in-1 (R-14), and drawing in perspective with the model built from construction lines
(R-15). The HCI reports in \code{docs/hci/} measure the editor with artists from Unity, Unreal and Adobe's tools in
mind, on keyboard, pen and tablet; the later rounds added the VRChat avatar and world pipelines, with the editor
reading and writing Unity's own project files.

\section{What this book covers}
Chapters~\ref{ch:context} to \ref{ch:bridges} describe the system: its context, the layers and the two flows that
organise everything (one frame, one edit), the scene and mesh model, rendering, the editor toolkit, and the bridges to
other programs. Chapter~\ref{ch:research} lists the research features and where each paper landed. Chapters~\ref{ch:sdlc}
to \ref{ch:performance} describe how the program is made: the agentic software development lifecycle with its
artifacts and gates, the build, test and release machinery, and the performance work. Chapter~\ref{ch:scope} says what
is not recreated. Chapter~\ref{ch:licensing} states the licence and the credits. The References chapter and the two
indexes close the book.

\section{Conventions}
\code{Monospace} is a file, a function, a flag or a console command; paths are relative to the \code{blendity/} folder
unless they start with \code{blender/}, which is Blender's source tree beside the project. \adr{NNNN} is a decision
record in \code{docs/decisions/}, \task{NNNN} a specification in \code{docs/tasks/}, and \rnd{NNN} an entry of
\code{CHANGELOG.md} (one entry per round of requests). Numbers quoted here were measured on the machines the CHANGELOG
names; where a figure is an estimate, the text says so.

\section{The numbers at 0.13.0}
\begin{center}
\begin{tabular}{lr}
\toprule
Source files under \code{src/} (C++, headers, Objective-C++, CUDA) & 380 \\
Lines in \code{src/} & about 297,000 \\
Lines of unit tests (\code{tests/}) & about 73,600 \\
Unit checks / tests, the build with Blender's libraries & 15,902 / 1,543 \\
Unit checks, the dependency-free build (Windows and Linux) & 15,253 / 15,259 \\
Test groups (\code{--group}) & 16 \\
Decision records & 104 \\
Task specifications & 379 \\
Rounds in the CHANGELOG & 409 \\
Learn lessons & 45 \\
Windows (module registry) & 26 \\
Keymap presets / editor feels & 7 / 4 \\
Commits on \code{main} since 2026-10-06 & 436 \\
\bottomrule
\end{tabular}
\end{center}
